\chapter{Background and Related Work}
\label{chap:background}

This chapter reviews the concepts and prior art on which the present work is
built. Section~\ref{sec:relational-vs-vector} contrasts relational and vector
databases as two complementary ways of organising data for retrieval.
Section~\ref{sec:llms} gives a short account of transformer-based language
models and the limitations that motivate retrieval augmentation.
Section~\ref{sec:embeddings} reviews the evolution of text representations up
to the dense sentence embeddings used in this work.
Section~\ref{sec:rag} describes the original Retrieval-Augmented Generation
architecture of Lewis et al.~\cite{Lewis2020RAG}. Finally,
Section~\ref{sec:related-work} positions the present work in the broader
landscape of RAG systems.

\section{Relational vs.\ Vector Databases}
\label{sec:relational-vs-vector}

Relational databases, introduced by Codd in the early 1970s, organise data in
tables with a fixed schema and allow exact retrieval through declarative SQL
queries. They excel at transactions, joins, and aggregations over structured
records, and modern engines such as PostgreSQL additionally offer full-text
search indexes built on top of inverted lists. Full-text search matches
\emph{lexical} tokens: a query for \texttt{contract OF employment} retrieves
documents that contain these exact words (possibly with stemming and
stop-word removal), but will miss a paraphrased passage such as
\texttt{labour agreement} even though the meaning is the same.

Vector databases take a different approach. Instead of indexing the tokens
that occur in each document, they index a dense, fixed-dimensional vector
that captures its \emph{semantic} content --- an \emph{embedding} --- computed
by a neural encoder. At query time, the same encoder produces an embedding
for the query, and the database returns the $k$ nearest neighbours in the
embedding space under a similarity metric, typically cosine similarity or
inner product. Because semantically similar passages are mapped to nearby
points even if they share no surface tokens, vector search naturally handles
paraphrases, synonyms, and translations.

Exact nearest-neighbour search is $\mathcal{O}(nd)$ in the number $n$ of
indexed vectors and their dimension $d$, which quickly becomes impractical.
Modern vector databases therefore rely on \emph{approximate nearest-neighbour}
(ANN) indexes that trade a controlled amount of recall for dramatic speed-ups.
Johnson et al.~\cite{Johnson2021FAISS} introduced FAISS, which combines
product quantisation with inverted-file indexes and GPU acceleration to
support billion-scale search on a single machine. More recent systems such as
HNSW (Hierarchical Navigable Small Worlds) have become the default choice in
practical vector databases, including ChromaDB, because they offer a
favourable trade-off between indexing cost and query time.

The two paradigms are not mutually exclusive. A typical production system
stores structured metadata (users, documents, audit logs) in a relational
database, and delegates semantic retrieval to a vector database. This is
exactly the architecture adopted by the present work, where PostgreSQL holds
the user accounts and ChromaDB holds the per-user collections of document
chunk embeddings.

\section{Large Language Models}
\label{sec:llms}

The transformer architecture introduced by Vaswani et
al.~\cite{Vaswani2017Attention} replaces the sequential recurrence of earlier
neural language models with a fully parallel self-attention mechanism. Each
layer computes, for every position in the sequence, a weighted combination of
all other positions, where the weights are derived from the compatibility
between pairs of positions. The resulting architecture scales efficiently
during both training and inference, and it has become the foundation of
virtually every state-of-the-art natural language processing model published
since 2018.

Large generative models trained on web-scale corpora --- the GPT family, the
Llama family, Google's Gemini family --- can produce coherent text over long
contexts and answer questions across an impressive range of topics. However,
they share three structural limitations that are directly relevant to this
work.

\textbf{Frozen knowledge.} The facts a model knows are encoded in its weights
during pre-training. Anything that happened after the training cutoff, or any
document that was not part of the pre-training corpus, is effectively
invisible to the model.

\textbf{Hallucination.} When asked about a topic the model has not seen, it
will still produce fluent output, often in the form of confident but
fabricated statements. This behaviour is particularly damaging in
domain-specific applications, where users cannot easily tell correct answers
from invented ones.

\textbf{Limited context window.} Even the largest generative models accept
only a bounded amount of text in a single request. They cannot simply be
given the entire corpus of a company at inference time. Some relevant subset
must be selected, and the selection is the job of a retrieval component.

These three limitations define the niche that RAG fills: retrieve the
relevant subset at query time, feed it to the model as context, and ground
the answer in auditable evidence.

\section{Text Embeddings}
\label{sec:embeddings}

The first generation of text representations was based on sparse
bag-of-words counts, weighted by TF-IDF. Such representations match queries
to documents by literal token overlap and fail on synonymy and paraphrase.
The introduction of word2vec and GloVe embeddings in the mid-2010s replaced
each word with a dense low-dimensional vector learnt from co-occurrence
statistics, capturing semantic regularities such as \emph{king - man +
woman $\approx$ queen}.

Word-level embeddings are a poor starting point for retrieval, however,
because they do not compose straightforwardly into sentence-level or
paragraph-level representations. The breakthrough came with contextual
encoders based on the transformer, BERT being the most influential. Reimers
and Gurevych~\cite{Reimers2019SentenceBERT} showed how to fine-tune BERT in a
Siamese configuration so that its \texttt{[CLS]} output becomes a sentence
embedding suitable for cosine-similarity retrieval; their Sentence-BERT
framework produces embeddings that remain the default baseline for many
semantic-search applications today. The same paper also demonstrated that
dense encoders can outperform classical retrieval methods on semantic
textual similarity benchmarks.

A parallel line of work focused on the retrieval task itself. Karpukhin et
al.~\cite{Karpukhin2020DPR} introduced \emph{Dense Passage Retrieval} (DPR),
a dual-encoder architecture in which one BERT encodes the query and a second
one encodes the passages, both trained end-to-end with a contrastive
objective that pulls positive pairs together and pushes negative pairs
apart. DPR outperformed the traditional BM25 baseline on open-domain
question answering and established dense retrieval as a serious competitor
to lexical methods.

For the practical system described in this thesis, we use Google's
\texttt{text-embedding-004} model through the Gemini API. This is a
pre-trained encoder that exposes a single embedding per input string and is
optimised for retrieval; we do not fine-tune it, because the scale and
licensing of the model makes fine-tuning both unnecessary and impractical
for the scope of a master's thesis.

\section{The Retrieval-Augmented Generation Paradigm}
\label{sec:rag}

Lewis et al.~\cite{Lewis2020RAG} formalised Retrieval-Augmented Generation
as an end-to-end architecture that jointly trains a dense retriever and a
generative decoder. Given an input query $x$, the retriever returns the
top-$k$ documents $z_1,\ldots,z_k$ from a non-parametric memory indexed as
dense vectors, and the generator produces an output $y$ conditioned on both
$x$ and the retrieved documents. The paper explored two variants ---
RAG-Sequence, which conditions on a single retrieved document per generated
token, and RAG-Token, which marginalises over the retrieved documents at
every step --- and demonstrated improvements over purely parametric
baselines on several open-domain QA tasks.

In the years since the original paper, the term \emph{RAG} has broadened to
denote any system that combines a retrieval step with a prompt-based
generation step, regardless of whether the retriever and the generator are
jointly trained. The practical recipe is now:

\begin{enumerate}
  \item split the corpus into chunks of a few hundred tokens;
  \item encode each chunk with a pre-trained encoder and store the vector
        in an ANN index;
  \item at query time, encode the question, retrieve the top-$k$ chunks,
        and build a prompt of the form
        \texttt{answer the question using the context below};
  \item pass the prompt to a frozen LLM and return the answer together with
        the citations derived from the chunk metadata.
\end{enumerate}

This recipe sacrifices some of the end-to-end elegance of the original paper
but dramatically simplifies deployment, because every component can be
swapped for a commercial API or an off-the-shelf vector store. The present
work follows this practical recipe.

\section{Related Work}
\label{sec:related-work}

Several open-source and commercial systems now offer RAG as a turnkey
product. Among the open-source libraries, LangChain and LlamaIndex provide
building blocks for document loading, chunking, embedding, retrieval, and
prompt templating, and integrate with a wide range of vector stores. Among
the commercial systems, Pinecone, Weaviate Cloud, and Google's Vertex AI
Search expose managed vector databases with integrated embedding and
retrieval pipelines.

The contribution of the present work is not to compete with these systems
on completeness, but to build a focused, well-engineered reference
implementation from first principles, strictly isolated per user, and
documented in enough detail to serve as a teaching example. The codebase
avoids heavy orchestration frameworks in order to keep every step visible:
a student reading the sources should be able to trace a question from the
HTTP handler all the way down to the ChromaDB similarity search and the
Gemini API call, without indirection.

The use of \emph{per-user collections} as the isolation mechanism, rather
than filter predicates on a shared collection, is deliberately simple and
trades some scalability at very large user counts for strong semantic
guarantees: it is impossible, by construction, for a query on one user's
collection to leak results from another.
